\section{Xây dựng hệ thống AI}

Hệ thống \textbf{InsightTune} có tích hợp hai dịch vụ (service) AI, hai dịch vụ này được triển khai sao cho tương thích tốt với toàn bộ
hệ thống và có tính đặc trưng cao (phân biệt với những dịch vụ AI cho đa tác vụ như: Gemini hay ChatGPT). Chúng tôi xin trình bày
những kiến thức nền tảng để xây dựng hiệu quả hai dịch vụ này qua hai phần riêng biệt dưới đây: dịch vụ gợi ý và trợ lý ảo.

\subsection{Dịch vụ gợi ý}
Để xây được một dịch vụ gợi ý tốt trong một thống lớn của \textbf{InsightTune}, ta cần hiểu hai miền kiến thức chính:
\begin{itemize}
    \item Một là về kiến trúc hệ thống, để từ đó hiểu được đầu vào và đầu ra của dịch vụ cũng như phương thức giao tiếp với những tác vụ khác của hệ thống.
    \item Hai là về xây dựng và triển khai một mô hình gợi ý.
\end{itemize}
Những yêu cầu kiến thức về kiến trúc hệ thống đã được trình bày tại \ref{sec:arch_knowledge} nên ở phần này chúng tôi chỉ tập trung vào giải thích ngắn
gọn những kiến thức xây dựng và triển khai mô hình gợi ý.

\subsubsection{Xây dựng}
Xây dựng mô hình gợi ý tức là chuẩn bị tập dữ liệu mẫu chia ra làm tập huấn luyện và tập kiểm tra, sau khi hoàn tất quá trình huấn luyện
trên tập huấn luyện, ta thu được một mô hình hoạt động tốt trên tập kiểm tra.

Kiến thức cần có để xây dựng lại một mô hình gợi ý như chúng tôi làm đó là:
\begin{itemize}
    \item Kiến thức về học tăng cường
    \item Xử lý dữ liệu
    \item Các thư viện học máy trong Python: Pytorch, numpy, pandas
    \item Công cụ đo lường và biểu diễn: wandb
\end{itemize}

\subsubsection{Triển khai}
Sau khi có được một mô hình gợi ý, chúng tôi triển khai mô hình như một dịch vụ trong hệ thống. Điều này phải tương thích với yêu cầu đầu vào và đầu ra của 
dịch vụ.

Những kiến thức cần có để triển khai hiệu quả là:
\begin{itemize}
    \item Kiến thức xây dựng dịch vụ bằng FastAPI
    \item Thiết lập đóng gọi dịch vụ bằng Docker
\end{itemize}

\subsection{Trợ lý ảo}
 
Chúng tôi xây dựng một hệ thống chatbot agent kết hợp với các tool:
\begin{itemize}
    \item Duckduckgo tool: là công cụ tra cứu, tìm kiếm thông tin trên mạng.
    \item Scrape Website tool: là công cụ để lấy dữ liệu từ một trang web.
\end{itemize}

Hệ thống sử dụng LLM là gemini-2.5-pro làm model chính để xử lí input đầu vào và kết hợp những tool để trả lời người dùng.
Dưới đây là kiến thức cần có để xây dựng hệ thống chatbot agent này:
\begin{itemize}
    \item Nguyên lý hoạt động của LLM và cơ chế function calling/tool use.
    \item Kiến thức Agent-base Chatbot.
    \item Thiết kế và tích hợp tools vào LLM.
    \item Kỹ năng viết prompt hiệu quả.
    \item Framework Langgraph để xây dựng stateful multi-agent workflows.
    \item Quản lí bộ nhớ hội thoại (memory management) với Checkpointer.
\end{itemize}

Và được triển khai thông qua FastAPI.

\subsection{Aider -- Công cụ lập trình cặp đôi với AI}
\label{sec:aider}

Aider~\cite{aider2024} là một công cụ lập trình cặp đôi (pair programming) mã nguồn mở chạy trên dòng lệnh (command-line), cho phép lập trình viên làm việc trực 
tiếp cùng với các mô hình ngôn ngữ lớn (LLM) để xây dựng và cải tiến mã nguồn. Công cụ này hỗ trợ hầu hết các LLM phổ biến hiện nay như GPT-4o, Claude 3.7 Sonnet, 
Gemini và DeepSeek, đồng thời tương thích với hơn 100 ngôn ngữ lập trình khác nhau.

\subsubsection{Đặc điểm nổi bật}

Aider có một số đặc điểm nổi bật giúp tăng năng suất phát triển phần mềm:
\begin{itemize}
    \item \textbf{Lập bản đồ mã nguồn (Repo map):} Aider tự động tạo bản đồ tổng quát của toàn bộ codebase, giúp LLM hiểu cấu trúc dự án và làm việc hiệu quả hơn 
    trong các dự án lớn, ngay cả khi chỉ có một phần nhỏ của mã nguồn được đưa vào ngữ cảnh.
    \item \textbf{Tích hợp Git:} Aider tự động tạo các commit với thông điệp có nghĩa sau mỗi thay đổi, giúp lập trình viên dễ dàng theo dõi lịch sử chỉnh sửa, 
    so sánh (diff) và hoàn tác (undo) khi cần.
    \item \textbf{Hỗ trợ đa LLM:} Aider kết nối được với hầu hết các nhà cung cấp LLM qua API, bao gồm cả các mô hình cục bộ (local models), đảm bảo tính 
    linh hoạt trong việc lựa chọn mô hình phù hợp.
    \item \textbf{Tích hợp IDE:} Aider có thể được sử dụng trực tiếp trong các trình soạn thảo mã nguồn thông qua tính năng theo dõi file (file watching), 
    cho phép lập trình viên thêm chú thích vào mã để kích hoạt Aider thực hiện thay đổi.
    \item \textbf{Hỗ trợ hình ảnh và trang web:} Aider cho phép đính kèm ảnh chụp màn hình, ảnh thiết kế hoặc trang web vào cuộc trò chuyện để cung cấp 
    ngữ cảnh trực quan cho LLM.
    \item \textbf{Giọng nói thành code (Voice-to-code):} Aider hỗ trợ nhập lệnh bằng giọng nói, giúp lập trình viên giao tiếp với AI mà không cần gõ phím.
    \item \textbf{Kiểm tra tự động (Lint \& Test):} Sau mỗi lần thay đổi, Aider tự động chạy linter và bộ kiểm thử để phát hiện lỗi và có thể tự khắc phục 
    những vấn đề được phát hiện.
\end{itemize}

\subsubsection{Ứng dụng trong dự án}

Trong quá trình phát triển hệ thống \textbf{InsightTune}, chúng tôi sử dụng Aider như một trợ lý lập trình AI để tăng tốc độ xây dựng và gỡ lỗi mã nguồn. 
Cụ thể, Aider được vận dụng trong các tác vụ như viết boilerplate code, tái cấu trúc (refactor) module, sinh tài liệu (docstring) và hỗ trợ debug. 
Việc sử dụng Aider giúp rút ngắn đáng kể thời gian phát triển, đồng thời đảm bảo chất lượng mã nguồn nhờ vào khả năng phân tích toàn diện codebase của công cụ này.